\section{Down-up scalars and ranks on the original member}

% BEGIN L86M1
\paragraph{L86M1. The level restriction of the theta operator.}
Fix a member $D=w.D(Z)$ admitted by the original layer point. Its
realization supplies the pinned characterisation, primitive conditions,
and the two readings $Y,Z_L$. Take $1\le j\le m$, write $j=i+1$ and
fix a used tuple $\mu$, either $\mu_0$ or $\mu_{\rm At}(w')$.
Choose $t=Y.\mathrm{tupT}(\mu)$, so $Y.\mathrm{tup}(t)=\mu$.
Write $V,W$ for the source and target Hilbert models at $i,t$ and
the distinguished embedding, $F=R.\mathrm{thG}(i,t):V\to W$, and
$\theta:M_t(\mu,j)\to M_j$ for the member's down lift.
Use the native coefficient embedding on both sides. If the place
coordinate uses its conjugate, transport both spaces, actions and
characters by the same conjugation. The pin gives $R.\mathrm{phs}(i)=1$ in the conjugated formula for the lift.
Inner-product identities and real eigenvalues survive this transport.

Put $P=P_0$, $N=|P|$, and $\xi(h)=\mathrm{emb}_{LL}(\psi_0(h)^{-1})$.
The subgroup acts unitarily on $V,W$. On $V$ the action extends the
action on rational forms, which preserves the class-sum pairing.
Those forms span $V$ by the Hilbert-model basis. On $W$ unitarity is
the member's primitive group condition. The primitive equivariance
identity for $F$ holds on the same rational spanning set, hence on
$V$. On either space put
\[
       E=\frac{1}{N}\sum_{h\in P}\xi(h)^{-1}\rho(h).
\]
Translation of the summation variable identifies its image with the
$\xi$-space. Unitarity, $|\xi(h)|=1$ and inversion in $P$ give
$E^*=E$. In particular $FE_V=E_WF$ and $E_VF^*=F^*E_W$.

There are maps $A:M_t(\mu,j)\to V$ and $B:M_j\to W$ with
\[
 \operatorname{im} A=E_V V,\qquad
 \operatorname{im} B=W_0\subset E_W W,\qquad
 A^*A=N\,1,\quad B^*B=N\,1,\qquad B\theta=FA.
\]
Indeed, the native source has the basis obtained from the rational
level columns by $j_{M_t}e_{M_t}$. Map them by the level inclusion
followed by $j_V$. The inclusion columns form a basis for the image
of the subgroup average. If $C$ is their matrix and $H$ the uncut
Hermitian form, the source metric is $N^{-1} C^*HC$. This proves
the image and metric assertions for $A$. On the target take
$B=Y.\mathrm{mj}.eW_i$. The target reading gives its squared-norm
factor $N$ and its image, the base-level space $W_0$; the
primitive level condition puts this space in the $\xi$-space.
Here $W_0$ also equals the span of the base-level target columns:
the auxiliary level equality identifies the model space with that
span, and scalar extension of the column basis gives the same range.
For the last identity it suffices to check a rational source column.
Its native and model readings evaluate the same theta integral
on pure tensors in the lowest compact types. The source realization
is injective, including when $t$ is replaced by its chosen tuple
representative. The difference of the two equivariant target
multiplicity maps vanishes on those tensors. Each lowest type
generates its irreducible discrete-series factor, so the difference
vanishes on the generated tensors and, by continuity, on the full
smooth model. The target evaluation is injective, which proves the
claimed equality on the column and hence everywhere.
% END L86M1

% BEGIN L86M2
\paragraph{L86M2. The restricted adjoint and the dimension of every lift image.}
For any $y\in M_j$ we have $E_WBy=By$. L86M1 gives
$F^*By\in E_VV=\operatorname{im} A$. For all $x\in M_t(\mu,j)$,
\[
 \langle Ax,F^*By\rangle
 =\langle B\theta x,By\rangle=N\langle\theta x,y\rangle
 =\langle Ax,A\theta^*y\rangle .
\]
Both vectors tested against $\operatorname{im}A$ lie in that space,
so $A\theta^*y=F^*By$. Therefore
$B\theta\theta^*y=FF^*By$. The native up lift is $\theta^*$.
For every parameter part $W_{ip}$, the primitive scalar identity
says $FF^*z=\lambda_{itp}z$ on $W_{ip}$, with
$\lambda_{itp}=R.\mathrm{lam3}(i,t,p,v_0)>0$. The primitive
target decomposition is
\[
 W_0=\bigoplus_p (W_{ip}\cap W_0).
\]
On the inverse image under $B$ of each summand,
$\theta\theta^*$ is multiplication by $\lambda_{itp}$, since $B$
is injective. The inverse on this finite sum multiplies its $p$-th
component by $\lambda_{itp}^{-1}$. It follows that $\theta^*$ is
injective and that its image has dimension $\dim M_j$, independently
of $\mu$. For a twisted place the tuple identity
$\mu_{\rm At}(w')\in\{\mu_0,\mu_{\rm Tw}(w')\}$ and the twisted-place
guard give $\mu_{\rm Tw}(w')=\mu_{\rm At}(w')$. Applying the image
dimension formula to it and to $\mu_0$ proves the original member's
rank comparison at every allowed $j$, for every native size datum.
% END L86M2

% BEGIN L86M3
\paragraph{L86M3. One family part for a finite type.}
Let $\sigma$ be a native finite type. Its reading is an isomorphism
class of irreducible smooth finite-adelic representations. Choose a
representative $\pi$; the conjugate row representation occurs at one
of the two row levels. At a good spherical place $v$ put $K$ equal
to its hyperspecial group. Each spherical double-coset operator
$T_d$ has a single scalar $c(v,d)$ on $\pi^K$. To see this, take a
nonzero vector in $\pi^K$. Its full orbit spans $\pi$ by
irreducibility. Averaging the orbit over $K$ shows that its $K$-fixed
space is generated by the spherical double-coset sums together with
translations at the complementary places. This holds for every nonzero
fixed vector. These operators thus act irreducibly on $\pi^K$ and
$T_d$ commutes with them, by commutativity of the unramified spherical
algebra and of the distinct local factors. The smooth cyclic
representation has countable dimension: an open stabilizer has
countably many translates in the second-countable finite-adelic
group. On a countable-dimensional simple complex module a commuting
operator $S$ is scalar. Indeed, if every $S-a$ is invertible then,
for fixed $u\ne0$, the uncountably many $(S-a)^{-1}u$ have a
nontrivial finite dependence. Clearing distinct denominators yields
a nonzero polynomial in $S$ killing $u$, although each of its linear
factors is invertible. Otherwise simplicity gives $S=a$ for some
$a$. Where the fixed space vanishes the scalar can be chosen
arbitrarily. At unused indices extend $c$ arbitrarily.

The good-level vectors of an isotypic space lie in the joint
eigenspace: each is a finite sum from equivariant images of $\pi$;
averaging its preimages over $K$ fixes the sum and makes each term
$K$-fixed. The pinned standard frames transfer the same
double-coset sums to the lift and row models. The row model uses
$\bar\pi$, giving eigenvalue $\overline{c(v,d)}$. Inversion gives
the adjoint double coset $d^\vee$ for the invariant inner product:
the two finite coset sums are convolution by the characteristic
functions of inverse double cosets, with the same normalisation.
On any nonzero occurrence this gives
$\overline{c(v,d)}=c(v,d^\vee)$, the convention in the defining
row eigenspace $E(c)$. Thus some nonzero row vector lies in
$E(c)$ at one of the two levels, and the lift reading of every
$\sigma$-vector at $j=i+1$ belongs to $W_i(c)$.

Write $p$ \emph{carries} $c$ if $E(c)\ne0$ and the adjoint of the
reference theta map takes $E(c)$ into the span of rational vectors of
inert type $p$. This condition depends only on $E(c)$. The $p$-th
row and lift parts are respectively the sums of $E(c')$ and
$W_i(c')$ for the systems $c'$ carried by $p$. If no $p$ carried
$c$, every nonzero such $E(c')$ would differ from $E(c)$ on some
tested double-coset eigenvalue. Otherwise the defining systems
of equations give the same subspace, which has the same carriers.
Distinct eigenvalues of one normal spherical operator give orthogonal
vectors; normality follows because its inverse-coset adjoint is in
the same commutative algebra. Our nonzero row vector would be
orthogonal to every row part. But at its own level the row space is
the sum of the intersections with these parts. This contradicts
positive definiteness. Choose a carrier $p=p(\sigma)$; for every
$i$ we have $B(\mathrm{isot}_j(\sigma))\subset W_{ip}$.
% END L86M3

% BEGIN L86M4
\paragraph{L86M4. The original field \textnormal{\texttt{lam}}.}
On a nonzero native isotypic summand L86M2 and L86M3 give
$\theta\theta^*y=\lambda_{it,p(\sigma)} y$. The member's reading
for $\mathrm{lamP}$ identifies its named real scalar with any real
scalar acting by this formula on that nonzero summand. It follows
that $\mathrm{lamP}(\mu,j,\sigma)=\lambda_{it,p(\sigma)}>0$.
On a zero summand the operator equation holds at its only vector
and the positivity implication has false antecedent. This covers
every used tuple and every allowed index. In particular a location
has $\mu=\mu_{\rm Tw}(w')=\mu_{\rm At}(w')$ by its twisted-place
guard, so the same calculation gives the location-wise field.
% END L86M4

% BEGIN L86M5
\paragraph{L86M5. Independence and the remaining dimension comparison.}
In the common ambient multiplicity representation the full
$\pi$-isotypic space is the algebraic sum of images of equivariant
maps out of $\pi$. Each nonzero image is simple. For any finite sum
of simple images one may discard an image contained in the preceding
sum; every other image meets that preceding sum trivially. The
resulting space is a finite direct sum of simple modules. Images of
a given simple module have zero projection to direct summands from
inequivalent simple modules, by Schur's lemma. Thus a finite
relation between sums belonging to pairwise inequivalent types has
each type contribution zero. The class reading distinguishes
different $\sigma$ by isomorphism class. Intersecting with the
level condition and using the injective target reading proves
independence of the native $\mathrm{isot}_j(\sigma)$.

Put $d_0=\dim E_{\rm base}$ for the reference base-level row space.
The member's dimension comparison equates
$\dim(W_{ip}\cap W_0)$ and $\dim(E_p\cap E_{\rm base})$ for every
$p$. The target and row direct sums and the isomorphism
$M_j\simeq W_0$ give $\dim M_j=d_0$. The base row reading identifies
$E_{\rm base}$ with the native space carrying the row basis. Consequently
$d_0=\sum_\sigma n_0(\sigma)$, where $n_0(\sigma)$ counts its basis
vectors of type $\sigma$. Independence now gives the exact remaining
test for the member's internal direct sum at row $j$:
\[
      \sum_\sigma\dim\mathrm{isot}_j(\sigma)=d_0 .
\]
It suffices to prove $\dim\mathrm{isot}_j(\sigma)\ge n_0(\sigma)$
for each type; summing and using independence forces equality. This
last inequality compares multiplicities of an individual
finite-adelic type across the two groups. More precisely, the upstream multiplicity input equates the cuspidal
multiplicity of an irreducible admissible good representation on
$Y_0$ with that of its transport by the finite-adelic isomorphism
to $W_i$, for the indicated archimedean discrete series. Once this
is supplied on the same member, the common finite level and good
projector act on the same finite representation, and taking the
level and projector subspaces preserves the multiplicity comparison.
The common spherical
label in L86M3 can contain multiple types and gives no such
lower bound by itself.
% END L86M5
